\subsection{Legal and Structural Lexer Constraints}

Before a Python name can become a runtime binding, it must survive the source-code reading phase. Python divides program text into tokens---names, keywords, literals, operators, delimiters, indentation markers---and rejects any assignment target that cannot be tokenized as a valid identifier. If the text fails here, execution never begins: no object is created, no name is bound, and no runtime type question exists yet.

This matters because Python names are flexible bindings, not C-style typed boxes. That flexibility starts only after the source code has been accepted. The syntax of the name is still fixed by the language at compile time.

The central lexical rule is:

\begin{quote}
A Python identifier must look like a valid name before it can be used as a runtime name.
\end{quote}

A valid ordinary identifier may contain letters, digits, and underscores, but it cannot start with a digit. It also cannot be one of Python's \emph{hard} reserved keywords---tokens such as \texttt{def}, \texttt{class}, and \texttt{import} that the parser reserves for grammatical structure.

However, hard keywords are not the only names Python treats specially. Built-in functions and types such as \texttt{print}, \texttt{id}, \texttt{type}, \texttt{len}, and \texttt{int} are \emph{not} protected compiler keywords. They are ordinary pre-populated entries in the \texttt{builtins} module dictionary. At runtime, a local or module-level assignment can mask them without triggering a syntax error:

\begin{verbatim}
print = 42          # legal source text; catastrophic runtime semantics
print("hello")      # TypeError: 'int' object is not callable
\end{verbatim}

The LEGB rule resolves \texttt{print} in the innermost namespace that currently defines it. If the local frame's \texttt{f\_locals} dictionary or the module's \texttt{f\_globals} dictionary contains \texttt{print = 42}, that binding wins over the built-in entry. CPython performs no compile-time warning because, from the compiler's perspective, \texttt{print} is simply another identifier-shaped name.

This is a deliberate design compromise. Python prioritizes namespace uniformity---every name is looked up through dictionary or fast-local machinery---over the C-like guarantee that certain identifiers are permanently reserved across all scopes. The execution risk is real: library code, interactive sessions, and hastily written modules can shadow built-ins and break downstream calls in ways that static analysis rarely catches.

Hard keywords remain the only names the parser refuses outright:

\begin{verbatim}
def = 42            # SyntaxError: invalid syntax
class = "Springfield"  # SyntaxError: invalid syntax
\end{verbatim}

The \texttt{keyword} module exposes the current interpreter's hard-keyword list (\texttt{keyword.kwlist}), which may grow between language versions. Soft keywords such as \texttt{match} and \texttt{case} occupy a middle tier: reserved only in specific grammatical positions. For identifier validation, the practical compile-time test remains: identifier-shaped \emph{and} not a hard keyword.

\subsubsection{Digit Restrictions: Why Identifiers Cannot Start with a Numerical Character}

A Python identifier cannot begin with a digit:

\begin{verbatim}
42answer = "value"   # SyntaxError: invalid decimal literal
\end{verbatim}

The rejection occurs at lex time, not at runtime. When Python reads source code, a leading digit begins a numeric literal. Allowing digit-leading identifiers would force the lexer to disambiguate forms such as \texttt{123abc}---integer literal followed by name, single identifier, or invalid number---on every token boundary. Python avoids this ambiguity with a clean rule: digits may appear after the first character, but not as the first character.

Valid forms include \texttt{answer42}, \texttt{fact\_7}, and \texttt{room101}. Invalid forms include \texttt{42answer}, \texttt{7fact}, and \texttt{101room}. The string method \texttt{isidentifier()} checks structural shape only; it does not create a variable and does not consult the keyword table.

Attempting to compile invalid source confirms the failure mode:

\begin{verbatim}
compile('42answer = 1', '<example>', 'exec')
# SyntaxError: invalid decimal literal
\end{verbatim}

The exact error message may vary between Python versions, but the structural rule is stable.

\subsubsection{Reserved Keywords: Why \texttt{def}, \texttt{class}, \texttt{import}, and Similar Tokens Cannot Be Used as Names}

Some words are reserved because Python uses them to express program structure. Examples include \texttt{def}, \texttt{class}, \texttt{import}, \texttt{if}, \texttt{else}, \texttt{for}, \texttt{while}, \texttt{return}, \texttt{try}, and \texttt{with}.

These tokens are forbidden not because they are poor English labels, but because they already have fixed grammatical jobs. If \texttt{def} were also an ordinary variable name, a line such as \texttt{def = 42} would make function-definition syntax structurally ambiguous. The parser must rely on \texttt{def} always introducing a function definition block.

The \texttt{keyword.iskeyword()} function distinguishes hard keywords from identifier-shaped words that happen to collide with built-ins:

%<<<PROGRAM:programs/leg01>>>
Hard keywords are rejected by the parser. Built-ins are not. That asymmetry is the architectural fault line between compile-time grammar and runtime namespace masking.

\subsubsection{Case Sensitivity: \texttt{name}, \texttt{Name}, and \texttt{NAME} as Distinct Bindings}

Python identifiers are case-sensitive. \texttt{name}, \texttt{Name}, and \texttt{NAME} are three distinct identifiers and therefore three independent bindings in whatever namespace defines them:

\begin{verbatim}
name = "Jerry"
Name = "George"
NAME = "Elaine"
# three separate dictionary entries; three separate objects
\end{verbatim}

Python matches spelling exactly, including letter case. The three names may even refer to objects of different types simultaneously. This is legal but architecturally hazardous: human readers scanning code often miss case-only distinctions, while the interpreter never does.

The complete compile-time rule for ordinary variable names:

\begin{quote}
A Python variable name must be a valid identifier, must not be a hard reserved keyword, and is matched by exact spelling including letter case. Runtime masking of built-ins remains possible even when all lexical rules are satisfied.
\end{quote}
